\documentclass[11pt,oneside]{article}

\usepackage[a4paper,margin=27mm,headheight=15pt]{geometry}
\usepackage[T1]{fontenc}
\usepackage[utf8]{inputenc}
\IfFileExists{libertinus.sty}{\usepackage{libertinus}}{\usepackage{lmodern}}
\usepackage{microtype}
\usepackage{amsmath,amssymb,amsthm,mathtools,mathrsfs}
\usepackage{bm}
\usepackage{booktabs,longtable,array,tabularx}
\usepackage{graphicx}
\usepackage{pgfplots}
\usepackage{xcolor}
\usepackage{enumitem}
\usepackage{fancyhdr}
\usepackage{titlesec}
\usepackage{listings}
\usepackage{xurl}
\usepackage{hyperref}
\usepackage[nameinlink,noabbrev]{cleveref}

\definecolor{linkblue}{RGB}{25,75,135}
\definecolor{shadegray}{RGB}{246,247,249}
\definecolor{rulegray}{RGB}{195,201,208}
\pgfplotsset{compat=1.18}
\hypersetup{
  colorlinks=true,
  linkcolor=linkblue,
  citecolor=linkblue,
  urlcolor=linkblue,
  pdftitle={The Fabius Function is Not Elementary},
  pdfsubject={Density of the analytic locus of an elementary function},
  pdfkeywords={Fabius function, elementary function, real analytic, nowhere analytic, isolated zeros, Lean, Mathlib}
}

\pagestyle{fancy}
\fancyhf{}
\fancyhead[L]{\small The Fabius function is not elementary}
\fancyhead[R]{\small\nouppercase{\leftmark}}
\fancyfoot[C]{\thepage}
\renewcommand{\headrulewidth}{0.4pt}

\titleformat{\section}{\Large\bfseries}{\thesection}{0.7em}{}
\titleformat{\subsection}{\large\bfseries}{\thesubsection}{0.7em}{}
\titleformat{\subsubsection}{\normalsize\bfseries}{\thesubsubsection}{0.7em}{}
\setcounter{secnumdepth}{3}
\setcounter{tocdepth}{2}
\setlist[itemize]{leftmargin=2em,itemsep=0.25em,topsep=0.4em}
\setlist[enumerate]{leftmargin=2.3em,itemsep=0.25em,topsep=0.4em}

\setlength{\emergencystretch}{2em}

\newtheorem{theorem}{Theorem}[section]
\newtheorem{proposition}[theorem]{Proposition}
\newtheorem{lemma}[theorem]{Lemma}
\newtheorem{corollary}[theorem]{Corollary}
\newtheorem{conjecture}[theorem]{Conjecture}
\theoremstyle{definition}
\newtheorem{definition}[theorem]{Definition}
\newtheorem{algorithm}[theorem]{Algorithm}
\newtheorem{example}[theorem]{Example}
\theoremstyle{remark}
\newtheorem{remark}[theorem]{Remark}
\newtheorem{warning}[theorem]{Warning}

\newcommand{\R}{\mathbb R}
\newcommand{\C}{\mathbb C}
\newcommand{\Q}{\mathbb Q}
\newcommand{\N}{\mathbb N}
\newcommand{\Z}{\mathbb Z}
\newcommand{\Up}{\operatorname{up}}
\newcommand{\e}{\mathrm e}
\newcommand{\ii}{\mathrm i}
\newcommand{\dd}{\,\mathrm d}
\newcommand{\E}{\mathbb E}
\newcommand{\Prob}{\mathbb P}
\newcommand{\bigO}{\mathcal O}
\newcommand{\supp}{\operatorname{supp}}
\newcommand{\Elem}{\mathcal E}
\newcommand{\An}{\mathcal A}
\newcommand{\repo}{\nolinkurl{Analysis/FabiusFunction}}

\newenvironment{proofidea}{\par\smallskip\noindent\textit{Proof architecture.}\ }{\hfill$\square$\par\smallskip}
\newenvironment{boxedremark}{%
  \par\smallskip\noindent\begin{center}\begin{minipage}{0.94\textwidth}%
  \color{black}\hrule\vspace{0.6em}\small
}{%
  \vspace{0.6em}\hrule\end{minipage}\end{center}\smallskip
}

\lstdefinestyle{pseudo}{
  basicstyle=\ttfamily\small,
  backgroundcolor=\color{shadegray},
  frame=single,
  rulecolor=\color{rulegray},
  columns=fullflexible,
  keepspaces=true,
  showstringspaces=false,
  xleftmargin=0.7em,
  xrightmargin=0.7em,
  aboveskip=0.8em,
  belowskip=0.8em
}

\title{\bfseries The Fabius Function is Not an Elementary Function\\[0.35em]
\large Density of the analytic locus, and what it costs to prove}
\author{Human-readable account of the formal development\\
\small Vladimir Reshetnikov's \texttt{ProveIt} repository}
\date{Repository snapshot: 25 August 2026}

\begin{document}
\maketitle

\begin{abstract}
The Fabius function $F$ is a $C^\infty$ probability distribution function on
$[0,1]$ that is real analytic at no point of that interval.  We prove, and
formalize in Lean~4 on top of \texttt{Mathlib}, that no elementary function of
one real variable agrees with $F$ on any nonempty open subset of $[0,1]$; in
particular the restriction of $F$ to $(0,1)$ is not elementary.

Since $F$ is $C^\infty$, no smoothness obstruction can decide the question:
the obstruction must be analyticity.  The whole content of the argument is
therefore a single positive statement about the class of elementary functions,
which we prove in the form we need it and no stronger:

\begin{center}
\emph{the set of points at which an elementary function is real analytic is a
dense open subset of $\R$.}
\end{center}

Density is exactly right and cannot be improved: $x\mapsto x^{-1}$ and
$x\mapsto |x|$ are elementary and fail to be analytic at the origin.  The
induction that proves it is short but has one genuinely delicate step, the
unary one, where a locally constant intermediate value must be separated from
a locally nonconstant one; the separation is furnished by the isolated-zeros
theorem for real analytic functions of one variable.
\end{abstract}

\tableofcontents

\section{Statement}

\subsection{The two inputs}

Write $\An(f)=\{x\in\R : f \text{ is real analytic at } x\}$ for the
\emph{analytic locus} of $f\colon\R\to\R$.  The proof combines one negative
fact about the Fabius function with one positive fact about elementary
functions.

\begin{theorem}[Nowhere analyticity; already formalized]\label{thm:nowhere}
Let $F$ be the Fabius function, extended by $0$ to the left of $0$ and by $1$
to the right of $1$.  Then $\An(F)=\R\setminus[0,1]$.  In particular $F$ is
real analytic at no point of $[0,1]$, endpoints included.
\end{theorem}

This is the unnumbered corollary to equation~(24) of Arias de Reyna,
\emph{Arithmetic of the Fabius function} (arXiv:1702.06487v3), transferred
from Rvachev's $\Up$ to $F$ and completed at the right endpoint; in the
repository it is \texttt{Fabius.canonical\_fabius\_analyticAt\_iff} in the
module \texttt{FabiusFunction.NowhereAnalytic}.  Nothing in the present
document reproves it.

\begin{theorem}[Density of the analytic locus]\label{thm:density}
If $f\colon\R\to\R$ is elementary, then $\An(f)$ is a dense open subset
of~$\R$.
\end{theorem}

\begin{corollary}[Main result]\label{cor:main}
Let $U\subseteq[0,1]$ be a nonempty open set and let $g\colon\R\to\R$ be
elementary.  Then $g$ and $F$ do not agree on $U$.  In particular there is no
elementary function that agrees with $F$ on $(0,1)$, and $F$ itself is not
elementary.
\end{corollary}

\begin{proof}
By \Cref{thm:density} the set $\An(g)$ is dense, so it meets the nonempty open
set $U$; pick $x\in U\cap\An(g)$.  If $g$ and $F$ agreed on $U$ they would
have the same germ at $x$, because $U$ is a neighbourhood of $x$; analyticity
at a point depends only on the germ, so $F$ would be analytic at $x\in[0,1]$,
contradicting \Cref{thm:nowhere}.
\end{proof}

\begin{boxedremark}
\Cref{cor:main} is the sharp form of the assertion ``the Fabius function is
not elementary'', and it is strictly stronger than the bare
non-elementarity of $F$ as a function on $\R$.  The bare statement leaves
open the possibility that some elementary function agrees with $F$ on all of
$(0,1)$ and differs from it outside; \Cref{cor:main} rules that out, and rules
it out on every subinterval.
\end{boxedremark}

\begin{remark}
Only a very weak consequence of \Cref{thm:density} is actually used: that an
elementary function is analytic at \emph{at least one} point of every nonempty
open set.  Density is nevertheless the statement one must prove, because it is
the statement that survives the induction: ``analytic somewhere'' is not
preserved by intersection, whereas ``dense open'' is.
\end{remark}

\section{The class of elementary functions}

\subsection{Definition}

We follow the description at
\url{https://en.wikipedia.org/wiki/Elementary_function}: elementary functions
are those obtained from polynomials, roots, exponentials, logarithms, and the
trigonometric functions and their inverses by finitely many applications of
the field operations and composition.  Formally:

\begin{definition}[Elementary functions]\label{def:elem}
The class $\Elem$ of \emph{elementary functions} $\R\to\R$ is the smallest
class of functions such that
\begin{enumerate}
\item every constant function $x\mapsto c$ lies in $\Elem$, and the identity
      $x\mapsto x$ lies in $\Elem$;
\item if $f,g\in\Elem$ then $x\mapsto f(x)+g(x)$, $x\mapsto f(x)g(x)$ and
      $x\mapsto -f(x)$ lie in $\Elem$;
\item if $f\in\Elem$ then $x\mapsto f(x)^{-1}$ lies in $\Elem$;
\item if $f\in\Elem$ and $r\in\R$ then $x\mapsto f(x)^r$ lies in $\Elem$;
\item if $f\in\Elem$ then each of
      $\exp\circ f$, $\log\circ f$, $\sin\circ f$, $\cos\circ f$,
      $\arcsin\circ f$, $\arctan\circ f$ lies in $\Elem$.
\end{enumerate}
\end{definition}

Two features of \Cref{def:elem} deserve comment.

\paragraph{Composition is a theorem, not an axiom.}
\Cref{def:elem} contains no clause ``$f,g\in\Elem\Rightarrow f\circ
g\in\Elem$''.  It does not need one.

\begin{proposition}\label{prop:comp}
$\Elem$ is closed under composition.
\end{proposition}

\begin{proof}
Induct on the derivation of the outer function $f$.  Every clause of
\Cref{def:elem} commutes with precomposition by a fixed $g$: the constant
$x\mapsto c$ composed with $g$ is again that constant; the identity composed
with $g$ is $g$ itself, which is elementary by hypothesis; and
$(f_1+f_2)\circ g=(f_1\circ g)+(f_2\circ g)$, $(\exp\circ f)\circ
g=\exp\circ(f\circ g)$, and likewise for every remaining clause.
\end{proof}

Consequently $\Elem$ contains every function one expects: differences
$f-g=f+(-g)$, quotients $f/g=f\cdot g^{-1}$, natural powers, all polynomial
and rational functions, $\sqrt{\cdot}$ and the $n$-th roots,
the absolute value $|x|=\sqrt{x^2}$, the tangent $\sin/\cos$, the hyperbolic
functions $\sinh x=(\e^x-\e^{-x})/2$, $\cosh$, $\tanh$, the arccosine
$\pi/2-\arcsin$, and the inverse hyperbolic sine
$\operatorname{arsinh}x=\log\bigl(x+\sqrt{1+x^2}\bigr)$.

\paragraph{Total functions and junk values.}
Every member of $\Elem$ is a \emph{total} function $\R\to\R$.  This is a
deliberate choice, and it is the choice that makes \Cref{cor:main} strongest.
The formalization inherits \texttt{Mathlib}'s conventions
\[
  0^{-1}=0,\qquad \log x=\log|x|,\qquad \log 0=0,\qquad
  0^{r}=0\ \ (r\ne 0),\qquad 0^{0}=1,
\]
\[
  x^{r}=|x|^{r}\cos(\pi r)\ \ (x<0),\qquad
  \arcsin x=\pm\tfrac{\pi}{2}\ \ (\pm x\ge 1),
\]
so that clauses (3)--(5) of \Cref{def:elem} never fail to produce a function.

\begin{boxedremark}
One might worry that junk values weaken the theorem, by admitting
``elementary'' functions that are not classically elementary.  They do the
opposite.  Admitting more functions into $\Elem$ makes the assertion ``$g\in
\Elem$ and $g=F$ on $U$'' easier to satisfy, hence its refutation
stronger.  And nothing is lost in the other direction: if a classical
elementary expression is well formed at every point of an open set $U$ --- no
vanishing denominator, no nonpositive logarithm, no real power of a
nonpositive base, no arcsine of a value outside $[-1,1]$ --- then on $U$ it
agrees pointwise with the total function produced by the same derivation,
since on $U$ the junk values are never consulted.  So every such classical
elementary function on $U$ is the restriction to $U$ of a member of $\Elem$,
and \Cref{cor:main} applies to it.
\end{boxedremark}

\begin{warning}[Odd roots need a different derivation]\label{warn:oddroot}
The condition ``no real power of a nonpositive base'' in the box above cannot
be dropped, and it is the one place where the total-function convention bites.
\texttt{Mathlib}'s $x^{r}$ at a negative base is $|x|^{r}\cos(\pi r)$, so
$(-8)^{1/3}=2\cos(\pi/3)=1$, whereas the classical real cube root of $-8$
is~$-2$.  The clause $x\mapsto f(x)^{r}$ of \Cref{def:elem} therefore supplies
the $n$-th roots only on $[0,\infty)$.

Nothing is lost, because the classical real odd root is in $\Elem$ by a
\emph{different} derivation,
\[
  x\ \longmapsto\ \frac{x}{|x|}\cdot|x|^{1/n},
\]
which uses only division, the absolute value and a real power, and which
evaluates to $0$ at $x=0$ under the same conventions.  This is
\texttt{IsElementary.signedRpow}, and
\texttt{Fabius.signedRoot\_pow} verifies formally that for odd $n$ its $n$-th
power is the identity, so it really is the real $n$-th root and not merely a
plausible substitute.
\end{warning}

\subsection{What is not claimed}\label{sec:notclaimed}

$\Elem$ contains the classical real odd-root functions but is not, as defined,
closed under passage to an arbitrary algebraic function: \Cref{def:elem} has no
clause producing a continuous branch of $P(x,y)=0$ for a polynomial $P$ with
elementary coefficients whose Galois group is not solvable.  Liouville's
differential field definition of ``elementary'' does admit such branches.  We
do not claim them here.  The obstruction is global formalization cost, not a
missing local theorem: \texttt{Mathlib}'s
\texttt{ContDiffAt.implicitFunction}, instantiated at analytic regularity
$\omega$, supplies the local analytic implicit-function theorem over $\R$.
What remains is a dense-analyticity theorem for continuous algebraic branches,
including a degree induction and control of the loci where the leading
coefficient or the derivative in the branch variable vanishes.

A second, much smaller gap concerns variable exponents.  \Cref{def:elem}
provides $x\mapsto f(x)^r$ for a \emph{fixed} real $r$, and
\texttt{IsElementary.rpow\_of\_pos} supplies $f^g$ whenever the base $f$ is
everywhere positive, which covers $x\mapsto x^x$ and every classical
variable-exponent expression.  The unrestricted two-variable power is not a
constructor, because \texttt{Mathlib}'s $x^y$ is sign-dependent at a negative
base --- there $x^y=|x|^y\cos(\pi y)$ --- and no single elementary formula
covers both signs at once.  Nothing is lost for \Cref{cor:main}: on any
interval on which the base has constant sign, $f^g$ does agree with a member
of $\Elem$, and \Cref{cor:main} already rules out agreement on every
subinterval.

We also make no claim about elementary functions of a complex variable, and
none about the closely related but distinct question of whether $F$ satisfies
an algebraic differential equation.

\section{Analyticity of the generators}

Each generator of \Cref{def:elem} is real analytic away from a finite set of
arguments.  All of these facts are in \texttt{Mathlib}, most of them in the
form ``$C^n$ for every $n$ in $\N\cup\{\infty,\omega\}$''; specializing to the
analytic exponent $\omega$ and applying \texttt{ContDiffAt.analyticAt}
converts them to the statements below.

\begin{lemma}\label{lem:generators}
\leavevmode
\begin{enumerate}
\item $\exp$, $\sin$, $\cos$ and $\arctan$ are real analytic at every point
      of $\R$.
\item $\log$ is real analytic at every $x\ne 0$ --- at negative arguments as
      well, because $\log$ is even under the convention $\log x=\log|x|$.
\item For each fixed $r\in\R$, the function $x\mapsto x^{r}$ is real analytic
      at every $x\ne 0$; at negative $x$ this uses the identity
      $x^{r}=|x|^{r}\cos(\pi r)$.
\item $\arcsin$ is real analytic at every $x\notin\{-1,1\}$; outside
      $[-1,1]$ it is locally constant.
\item $x\mapsto x^{-1}$ is real analytic at every $x\ne 0$.
\end{enumerate}
\end{lemma}

\begin{remark}
The bad sets are sets of \emph{values} of the inner function, not sets of
points of the domain --- $\{0\}$ for the reciprocal, the logarithm and a real
power, $\{-1,1\}$ for the arcsine, and $\varnothing$ for the rest.  This is
the shape the induction consumes, and it is why \Cref{lem:key} below is stated
with a finite set $B\subseteq\R$ of forbidden values.
\end{remark}

\section{Density of the analytic locus}

\subsection{Openness}

\begin{lemma}\label{lem:open}
For every $f\colon\R\to\R$ the set $\An(f)$ is open.
\end{lemma}

\begin{proof}
Analyticity at $x$ means that $f$ agrees near $x$ with the sum of a
convergent power series; that series then converges on a whole ball around
$x$ and, after a change of origin, represents $f$ near every point of that
ball.  So $\An(f)$ contains a neighbourhood of each of its points.
\end{proof}

\subsection{The unary step}

The following lemma is where the argument is actually made.  Recall the shape
of the difficulty.  If $f$ is analytic on a dense open set and $g$ is analytic
everywhere except at finitely many values, one would like to say that
$g\circ f$ is analytic wherever $f$ is analytic and $f(x)$ avoids the bad
values.  That set need not be dense: consider $f\equiv 0$ and $g(y)=y^{-1}$.
What saves the statement is that in that degenerate case $g\circ f$ is
constant, hence analytic anyway.  Separating the two cases is exactly the
isolated-zeros theorem.

\begin{lemma}[Key lemma]\label{lem:key}
Let $f,g\colon\R\to\R$ and let $B\subseteq\R$ be finite.  Suppose $\An(f)$ is
dense and $g$ is analytic at every $y\notin B$.  Then $\An(g\circ f)$ is
dense.
\end{lemma}

\begin{proof}
Let $U$ be a nonempty open set; we must produce a point of $U\cap\An(g\circ
f)$.  Since $\An(f)$ is dense there is $x_0\in U\cap\An(f)$, and by
\Cref{lem:open} the set $V=U\cap\An(f)$ is an open neighbourhood of $x_0$.
Distinguish two cases.

\emph{Case 1: $f$ is constant equal to some $c\in B$ on a neighbourhood of
$x_0$.}  Then $g\circ f$ equals the constant $g(c)$ on that neighbourhood, so
$g\circ f$ is analytic at $x_0$ itself, and $x_0\in U$.

\emph{Case 2: for no $c\in B$ is $f$ eventually equal to $c$ near $x_0$.}  Fix
$c\in B$.  The function $f-c$ is analytic at $x_0$, so by the isolated-zeros
theorem for analytic functions of one real variable, either $f-c$ vanishes
identically near $x_0$ --- excluded in this case --- or $f-c$ is nonvanishing
on some punctured neighbourhood of $x_0$.  Thus for each of the finitely many
$c\in B$ there is a punctured neighbourhood of $x_0$ on which $f\ne c$.
Intersecting these finitely many punctured neighbourhoods with the punctured
neighbourhood $V\setminus\{x_0\}$ leaves a punctured neighbourhood of $x_0$,
which is nonempty because $\R$ has no isolated points.  Any point $t$ of it
satisfies $t\in U$, $f$ analytic at $t$, and $f(t)\notin B$; hence $g$ is
analytic at $f(t)$ and $g\circ f$ is analytic at $t$.
\end{proof}

\begin{remark}
The case distinction is not an artefact of the write-up.  Without it the lemma
is false as stated, for the naive reason above.  What makes the two cases
exhaustive is the isolated-zeros theorem, applied to $f-c$: an analytic germ
either vanishes identically or vanishes nowhere on some punctured
neighbourhood.  A merely $C^\infty$ inner function admits neither
alternative --- it can equal $c$ on a sequence accumulating at $x_0$ without
being constant there --- so analyticity of $f$, and not just smoothness, is
what the step consumes.  This is the only place in the proof where a theorem
of substance about real analytic functions is used.
\end{remark}

\subsection{The induction}

\begin{proof}[Proof of \Cref{thm:density}]
Openness of $\An(f)$ is \Cref{lem:open} and holds for every $f$.  Density is
proved by induction on the derivation of $f$ in \Cref{def:elem}.

\emph{Constants and the identity.}  Here $\An(f)=\R$.

\emph{Sums and products.}  If $\An(f)$ and $\An(g)$ are dense then so is
$\An(f)\cap\An(g)$, because a finite intersection of dense open sets is dense.
Analyticity is preserved by sums and products pointwise, so
$\An(f)\cap\An(g)\subseteq\An(f+g)$ and likewise for $fg$; a superset of a
dense set is dense.

\emph{Negation.}  $\An(f)\subseteq\An(-f)$.

\emph{The unary clauses.}  Apply \Cref{lem:key} with the pair $(f,g)$ and the
finite bad set $B$ read off from \Cref{lem:generators}: $B=\varnothing$ for
$\exp$, $\sin$, $\cos$, $\arctan$; $B=\{0\}$ for $y\mapsto y^{-1}$, for
$\log$, and for $y\mapsto y^{r}$; and $B=\{-1,1\}$ for $\arcsin$.
\end{proof}

\begin{remark}[Sharpness]
Density cannot be strengthened to ``$\An(f)=\R$'': the elementary functions
$x\mapsto x^{-1}$, $x\mapsto|x|$ and $x\mapsto\arcsin x$ have analytic loci
$\R\setminus\{0\}$, $\R\setminus\{0\}$ and $\R\setminus\{-1,1\}$.  Nor can
``dense open'' be strengthened to ``open with finite complement'': the
elementary function $x\mapsto (\sin x)^{-1}$ omits the infinite discrete set
$\pi\Z$.  What density does give for free is that the complement is closed
and nowhere dense, being the complement of a dense open set.  Whether the
complement of the analytic locus of an elementary function is always
\emph{discrete} is a natural further question that we neither prove nor
use.
\end{remark}

\section{The formal development}

\subsection{Correspondence with the Lean names}

All statements below live in the namespace \texttt{Fabius}.

\begin{center}
\small
\begin{tabular}{@{}p{0.40\textwidth}p{0.31\textwidth}p{0.23\textwidth}@{}}
\toprule
\textbf{Object or statement} & \textbf{Lean name} & \textbf{Module} \\
\midrule
Elementary functions (\Cref{def:elem})
  & \texttt{IsElementary}
  & \texttt{ElementaryFunction} \\
Closure under composition (\Cref{prop:comp})
  & \texttt{IsElementary.comp}
  & \texttt{ElementaryFunction} \\
Derived closures: $-$, $/$, $x^n$, $\sqrt{\ }$, $|\cdot|$, $\tan$, $\sinh$, $\cosh$, $\tanh$, $\arccos$, $\operatorname{arsinh}$, polynomials
  & \texttt{IsElementary.sub}, \texttt{.div}, \texttt{.npow}, \texttt{.sqrt}, \texttt{.abs}, \texttt{.tan}, \texttt{.sinh}, \texttt{.cosh}, \texttt{.tanh}, \texttt{.arccos}, \texttt{.arsinh}, \texttt{.rpow\_of\_pos}, \texttt{.polynomial}
  & \texttt{ElementaryFunction} \\
Analytic locus $\An(f)$
  & \texttt{analyticLocus}
  & \texttt{ElementaryFunction} \\
Openness (\Cref{lem:open})
  & \texttt{isOpen\_analyticLocus}
  & \texttt{ElementaryFunction} \\
Generators (\Cref{lem:generators})
  & \texttt{analyticAt\_log\_of\_ne\_zero}, \texttt{analyticAt\_rpow\_const}, \texttt{analyticAt\_arctan}, \texttt{analyticAt\_arcsin}; for $x\mapsto x^{-1}$, $\exp$, $\sin$, $\cos$ the \texttt{Mathlib} lemmas \texttt{analyticAt\_inv}, \texttt{analyticAt\_rexp}, \texttt{Real.analyticAt\_sin}, \texttt{Real.analyticAt\_cos} are used directly
  & \texttt{ElementaryFunction} \\
Auxiliary: analyticity of $\sqrt{\cdot}$ off $0$ (not used in the induction, which routes $\sqrt{\cdot}$ through \texttt{rpow})
  & \texttt{analyticAt\_sqrt\_of\_ne\_zero}
  & \texttt{ElementaryFunction} \\
Signed root, and its defining property (\Cref{warn:oddroot})
  & \texttt{IsElementary.signedRpow}, \texttt{signedRoot\_pow}
  & \texttt{ElementaryFunction} \\
Membership of the named functions
  & \texttt{isElementary\_exp}, \texttt{\_log}, \texttt{\_sin}, \texttt{\_cos}, \texttt{\_tan}, \texttt{\_arcsin}, \texttt{\_arccos}, \texttt{\_arctan}, \texttt{\_sinh}, \texttt{\_cosh}, \texttt{\_tanh}, \texttt{\_arsinh}, \texttt{\_sqrt}, \texttt{\_abs}, \texttt{\_inv}, \texttt{\_rpow}, \texttt{\_id}
  & \texttt{ElementaryFunction} \\
Key lemma (\Cref{lem:key})
  & \texttt{dense\_analyticLocus\_comp}
  & \texttt{ElementaryFunction} \\
Density (\Cref{thm:density})
  & \texttt{IsElementary.dense\_analyticLocus}
  & \texttt{ElementaryFunction} \\
Analytic at a point of any nonempty open set
  & \texttt{IsElementary.exists\_analyticAt\_of\_isOpen}
  & \texttt{ElementaryFunction} \\
Exceptional set closed with empty interior
  & \texttt{isClosed\_compl\_analyticLocus}, \texttt{IsElementary.interior\_compl\_analyticLocus}
  & \texttt{ElementaryFunction} \\
Nowhere analyticity (\Cref{thm:nowhere})
  & \texttt{canonical\_fabius\_not\_analyticAt}, \texttt{canonical\_fabius\_analyticAt\_iff}
  & \texttt{NowhereAnalytic} \\
Main result (\Cref{cor:main})
  & \texttt{canonical\_fabius\_not\_isElementary\_on\_Ioo}
  & \texttt{NotElementary} \\
$F$ is not elementary
  & \texttt{canonical\_fabius\_not\_isElementary}
  & \texttt{NotElementary} \\
General open subset of $[0,1]$
  & \texttt{not\_isElementary\_eqOn}
  & \texttt{NotElementary} \\
Rvachev's $\Up$ and the signed extension
  & \texttt{rvachev\_not\_isElementary}, \texttt{extendedFabius\_not\_isElementary}
  & \texttt{NotElementary} \\
\bottomrule
\end{tabular}
\end{center}

All modules are under \texttt{\repo/Lean/FabiusFunction/}.

\subsection{Ingredients taken from \texttt{Mathlib}}

The formalization is short because \texttt{Mathlib} already contains the two
theorems of substance that the argument consumes.

\begin{itemize}
\item \texttt{AnalyticAt.eventually\_eq\_zero\_or\_eventually\_ne\_zero} and
      its two-function form: the isolated-zeros theorem for a function of one
      variable over a nontrivially normed field, which is what powers
      \Cref{lem:key}.  It applies verbatim to $f\colon\R\to\R$.
\item \texttt{isOpen\_analyticAt}: the analytic locus is open, which is
      \Cref{lem:open} verbatim.  (The pointwise form, analyticity propagating
      from a point to a neighbourhood, is
      \texttt{AnalyticAt.eventually\_analyticAt}; the development uses the
      set-level statement.)
\item \texttt{ContDiffAt.analyticAt} together with the family of
      \texttt{contDiffAt\_\ldots} lemmas stated for exponents in
      $\N\cup\{\infty,\omega\}$: this is how \Cref{lem:generators} is
      obtained, including the two facts that are otherwise awkward over $\R$,
      namely analyticity of $\arctan$ everywhere and of $\arcsin$ off
      $\{\pm1\}$.
\end{itemize}

\subsection{Statement of the main theorem in Lean}

\begin{lstlisting}[style=pseudo]
theorem canonical_fabius_not_isElementary_on_Ioo :
    not (exists g : R -> R, IsElementary g and
         Set.EqOn g (fabiusReal fabius) (Set.Ioo 0 1))
\end{lstlisting}

\noindent (rendered here in ASCII; the source uses the usual notation).  The
development contains no \texttt{sorry}, no \texttt{axiom} and no
\texttt{opaque}, and its axiom set is the \texttt{Mathlib} default
\texttt{propext}, \texttt{Classical.choice}, \texttt{Quot.sound}.

\section{Discussion}

\subsection{Why smoothness cannot decide the question}

The Fabius function is $C^\infty$ on $\R$, and its derivatives satisfy the
exact self-similar relation $F'(x)=2F(2x)$ on $[0,1/2]$.  Any proof of
non-elementarity must therefore distinguish $F$ from the elementary functions
by a property strictly finer than differentiability of every order.  Real
analyticity is the natural candidate, and \Cref{thm:density} says it is a
sufficient one.  It is also, in a precise sense, the \emph{only} such
candidate we know how to extract from \Cref{def:elem} directly.  Nothing in
this document rules out a different separating property; what it does show is
that analyticity suffices, and that it is available at the modest cost of
\Cref{lem:key}.

\subsection{The role of the endpoints}

\Cref{thm:nowhere} includes the endpoints $0$ and $1$, where $F$ is flat to
infinite order.  It is worth recording that \Cref{cor:main} does \emph{not}
need them.  A set that is open in $\R$ and contained in $[0,1]$ is
automatically contained in $(0,1)$ --- an $\R$-open set containing $0$
contains negative numbers --- so the witness produced by
\Cref{thm:density} is always an interior point, and the endpoint case of
\Cref{thm:nowhere} is never invoked.  The same is true if one reads ``open''
in the relative topology of $[0,1]$: a nonempty relatively open subset still
contains a nonempty $\R$-open subinterval of $(0,1)$.

What the endpoints do buy is exactness.  $\An(F)$ is precisely
$\R\setminus[0,1]$: outside the unit interval $F$ is locally constant, hence
analytic, so the nowhere-analytic locus is as large as it can be and no
larger.

\subsection{What a stronger theorem would require}

Two natural strengthenings are visible from here.

\begin{enumerate}
\item \emph{Liouvillian functions.}  Replacing $\Elem$ by the class of
      functions lying in an \emph{elementary} differential-field extension of
      the rational functions in Liouville's sense --- allowing arbitrary
      algebraic branches alongside exponentials and logarithms --- would
      subsume \Cref{def:elem}.  (The wider \emph{Liouvillian} class also
      admits antiderivatives, and so contains $\operatorname{erf}$,
      $\operatorname{Ei}$ and $\operatorname{Si}$, which are not elementary;
      that class is not what is meant here.)  The local analytic
      implicit-function theorem is already available.  The missing ingredient
      is the global dense-analyticity argument for algebraic branches, including its
      degree induction and singular-locus analysis, as recorded in
      \Cref{sec:notclaimed}.
\item \emph{Differential algebraicity.}  A stronger and genuinely different
      statement is that $F$ satisfies no algebraic differential equation
      $P\bigl(x,F,F',\ldots,F^{(n)}\bigr)=0$ with polynomial $P$.  Such a
      statement does not follow from nowhere analyticity --- differentially
      algebraic functions need not be analytic --- and is not addressed here.
\end{enumerate}

\end{document}
